\documentclass[10pt,a4paper]{article} % Reduced font size to 10pt
\usepackage{geometry}
\geometry{left=1in, right=1in, top=1in, bottom=1in}
\usepackage{enumitem}
\usepackage{hyperref}
\usepackage{titlesec}
\usepackage{xcolor}
\usepackage{kotex}
% Title formatting
\titleformat{\section}{\large\bfseries}{}{0em}{}[\titlerule]

% Custom commands
\newcommand{\seclabel}[1]{\section*{\uppercase{#1}}}
\newcommand{\cvitem}[2]{\textbf{#1} \hfill \small{#2}\\}

\hypersetup{
    colorlinks=true,
    linkcolor=cyan,
    filecolor=cyan,
    urlcolor=cyan,
    citecolor=cyan
}


\begin{document}

\begin{center}
    \textbf{\Huge{Kyojun Choo}} \\ \vspace{1.0em}
    B.S student \\ Soongsil University \\ School of Computer Science \& Engineering \\
    \textbf{Email: }\href{mailto:kjun5644@soongsil.ac.kr}{kjun5644@soongsil.ac.kr} \\
    \textbf{Home: }\url{https://chubeloop.github.io/} \\
    \textbf{LinkedIn: }\url{https://www.linkedin.com/in/kyojun-choo-8047b339a/} \\
    \textbf{Lab: }\url{https://nlp.ssu.ac.kr/}
\end{center}

\vspace{1em}

\seclabel{Lifelong Research Objective}
\small{My lifelong research objective is to build AI systems that people can use with confidence and trust. As AI becomes an increasingly important part of everyday life, I aim to ensure that its capabilities, limitations, and behaviors can be reliably understood and evaluated. Through research on model interpretability and trustworthy evaluation, I hope to contribute to a future where people can benefit from AI while clearly understanding when and why its outputs can be trusted.}

\seclabel{Research Interests}
\small{My research interests center on understanding and improving Large Language Models (LLMs) through interpretability and evaluation. I am particularly interested in analyzing how linguistic, task-related, and capability-related information is represented within model internals, and in using these insights to better explain model behaviors, limitations, and failure modes. At the same time, I aim to develop more reliable and meaningful evaluation methodologies that go beyond surface-level benchmark scores by examining robustness, validity, generalization, and the extent to which benchmarks truly reflect the capabilities they are intended to measure. Ultimately, I seek to connect interpretability with evaluation so that a deeper understanding of model representations can inform the design of more trustworthy, transparent, and practically useful LLM systems.}

\seclabel{Education}
\begin{itemize}[left=0pt]
    \item \cvitem{Soongsil University}{\small{Feb. 2021 - Present}}
    \small{B.S student School of Computer Science \& Engineering (advisor: Chanjun Park) \\
    \small{GPA: 4.14/4.5}}
\end{itemize}

\seclabel{Work Experiences}
\begin{itemize}[left=0pt]
    \item \cvitem{Natural Language Processing Lab, Soongsil University}{\small{Nov. 2025 - Present}}\small{Undergraduate Researcher}
    \item \cvitem{SH\&M Coding Class, Spartan Future Education Institute}{\small{March. 2025 - Present}}\small{Coding Instructor}
\end{itemize}

\newpage


\seclabel{Publications}
\textit{Equal contribution}: *, \textit{Corresponding author}: $\dagger$

\vspace{0.8em}

\input{publications.generated.tex}

\vspace{1.0em}


\section{Research Projects}
\textbf{Soongsil University}
\begin{itemize}[left=0pt]
    \item \small 2026년 인공지능 말평 운영 및 연계 행사 개최 (2026, 참여연구자 (용역과제), National Institute of the Korean Language)
    \item \small 2026년 인공지능의 한국어 능력 표준 평가 과제 구축 및 활용 (2026, 참여연구자 (용역과제), National Institute of the Korean Language)
    \item \small 우수연구 신진연구 (유형B) (2026-2029, 참여연구자 (3책5공), NRF)
    \item \small Foundation Model 운용과정에서 민감정보 추론 방지 및 리스크 평가 기술 개발 (2026-2028, 참여연구자 (3책5공), Korea Internet \& Security Agency (한국인터넷진흥원))
    \item \small 토론과 소통으로 지식을 연결하는 AI 에이전트 기술 개발 (2026-2028, 참여연구자\textbf{(PM)} (3책5공), Korea Creative Content Agency(한국콘텐츠진흥원))
    \item \small AI 보안 특화 'AI 에이전트' 개발 선행 연구 (2026,  참여연구자 (용역과제), 한국정보보호학회 산하 AI보안연구회)
    \item \small 국가·공공기관 가이드라인 준수를 위한 AI 시스템 보안 위협 자동화 기술점검 도구 개발 (2026, 참여연구자, KIST)
    \item \small AI 연구용 컴퓨팅 지원 프로젝트 (파라미터 수준의 정보 제어 원천기술 연구 및 파라미터 스튜디오 플랫폼 개발) (2026,  참여연구자, KETI)
\end{itemize}



\seclabel{Awards/Fellowships/Research Grants}


\noindent\textbf{Awards}

\begin{itemize}[left=0pt]
    \item \small{
    \textbf{Best Paper Award (HCLT 2026)} \\
    오답의 의미적 타당성 제거를 통한 LLM 벤치마크 결함 문항 식별 \\
    송민수, 김찬우, \underline{추교준}, 송은서, 장지윤, 남채린, 이시현, 박찬준\textsuperscript{$\dagger$} \\
    \textit{한글 및 한국어 정보처리 학술대회 (HCLT) 2026}
    }
\end{itemize}
\vspace{1.0 em}

\noindent\textbf{Fellowships}
\begin{itemize}
    \item \small Baekma Academic Excellence Scholarship (Full Tuition), Soongsil University (Mar. 2025)
    \item \small Baekma Academic Excellence Scholarship (Full Tuition), Soongsil University (Sep. 2021)
\end{itemize}

\noindent\textbf{Research Grants}
\begin{itemize}
    \item \small VESSL AI GPU Grant Program (2026)
\end{itemize}

\seclabel{Patents}

\noindent\textbf{Domestic Patent Applications}
\begin{itemize}[left=0pt]
    \item \textbf{Method and Apparatus for Determining Functional Requirements of User Queries Using Pre-Generation Hidden States of a Large Language Model} \\
    \textit{대규모 언어 모델의 생성 전 은닉 상태를 이용한 사용자 질의의 기능적 요구사항 판단 방법 및 장치} \\
    Chanjun Park, \underline{Kyojun Choo} \\
    Korean Patent Application, filed on August 19, 2026 (Application No. 10-2026-0156645).

    \item \textbf{APPARATUS AND METHOD FOR AUDITING PERFORMANCE RANKINGS AMONG ARTIFICIAL INTELLIGENCE MODELS USING BENCHMARK-SPECIFIC AUDIT CONDITIONS} \\
    \textit{벤치마크별 감사 조건을 이용한 인공지능 모델 간 성능 순위 감사 장치 및 방법} \\
    Chanjun Park, \underline{Kyojun Choo} \\
    Korean Patent Application, filed on Oct 6, 2026 (Application No. 10-2026-0190777).
\end{itemize}

\section{External Activities}

\textbf{Teaching Assistant}
\begin{itemize}[left=0pt]
    \item \small{SaRangAXConnect, AI Transformation (AX) Program (September 2026)}
\end{itemize}
\vspace{0.4em}


\end{document}
